\section*{Chapter \ref{ch:m3:power-markets-design} --- Power Markets I: Design}

\begin{solution}{exo:m3:power-markets-design:1}
\qty{40}{EUR/MWh}: the 300 MWh at zero and 150 of the 200 MWh at 40 are accepted, the
offer at 40 being marginal; the offers at 90 and 150 are rejected.
\end{solution}

\begin{solution}{exo:m3:power-markets-design:2}
$(30 + 0.202 \times 70)/0.5 = \qty{88.28}{EUR/MWh}$.
\end{solution}

\begin{solution}{exo:m3:power-markets-design:3}
In a uniform-price auction all accepted energy is worth the same at the margin, and
sellers can bid their marginal cost: any accepted offer earns the clearing price.
Under pay-as-bid each seller would bid what it expects the marginal price to be, not its
cost, so the offers would no longer reveal the \omterm{def:m3:power-markets-design:merit}{merit order}; the auction format itself is
studied in One Quant Book 4, chapter 29.
\end{solution}

\begin{solution}{exo:m3:power-markets-design:4}
Flow 3~GW from A to B, prices \qty{82.38}{} and \qty{89.68}{EUR/MWh}, rent $7.30 \times
3\,000 \approx$ EUR~21\,900 for the hour (EUR~21\,890 with the unrounded price gap,
\qty{7.297}{EUR/MWh}). Zone A has 63~GW of capacity cheaper than
\qty{89.35}{} against 50~GW of demand: with more than 13~GW of line, its gas plants can
export too and the prices converge at \qty{89.35}{EUR/MWh}.
\end{solution}

\begin{solution}{exo:m3:power-markets-design:5}
It produces while the price plus the premium is positive: it stops below
$-\qty{60}{EUR/MWh}$, that is in the four hours from 12:00 to 16:00 on 11 May 2025
(prices from $-212.82$ to $-110.06$).
\end{solution}

\begin{solution}{exo:m3:power-markets-design:6}
Inside a zone the network still has limits; the nodal price, which includes congestion,
differs across them, but the zonal market ignores this and clears at one price. The
operator then redispatches, paying plants on one side of the bottleneck to reduce and on
the other to increase output, and the cost is recovered through network charges from
all consumers.
\end{solution}

\begin{solution}{exo:m3:power-markets-design:7}
457 in 2024 and 573 in 2025; the most at 13:00 local time (171 over the two years); the
lowest price $-\qty{250.32}{EUR/MWh}$ at 13:00 on 11 May 2025.
\end{solution}

\begin{solution}{exo:m3:power-markets-design:8}
A negative price says that at that hour one more MWh has negative value: supply that cannot
or will not stop exceeds demand. It pays flexible consumers to consume and flexible producers
to stop, which is the market working; producers that keep producing do so because stopping
costs more (restarts) or because they are paid per MWh produced (subsidies).
\end{solution}

\begin{solution}{pb:m3:power-markets-design:1}
\textbf{1.} Nine hours, average $-\qty{98.54}{EUR/MWh}$. \textbf{2.} $-\qty{250.32}{}$, 13:00
to 14:00. \textbf{3.} From 11:00 to 16:00 (five hours). \textbf{4.} Exports to neighbouring
zones up to the lines' capacity, storage (pumped hydro, batteries), and curtailment of
producers that stop at negative prices. \textbf{5.} Supply bid below zero, from producers who
lose more by stopping or are paid per MWh produced, exceeded demand at zero.
\textbf{6.} $400 \times (82.38 + 250.32) = $ EUR~133\,080. \textbf{7.} $400 \times \sum_h
(82.38 - p_h) = $ EUR~651\,308 over the nine hours. \textbf{8.} Switching off avoids that
loss but costs a restart (fuel to warm up, wear, and the risk of not being ready for the
evening peak). \textbf{9.} Stay on if the restart cost exceeds the loss of running at minimum
load through the negative hours. \textbf{10.} A block offer over the day at minimum load with
a limit that includes the restart cost, or offers at negative prices for the minimum load.
\textbf{11.} Every negative hour whose price is above $-60$: 9:00, 10:00, 11:00, 16:00 and
17:00. \textbf{12.} About $-\qty{60}{EUR/MWh}$: their offers sit at minus their premium.
\textbf{13.} By paying no premium when prices are negative: in Germany, for new solar
plants since February 2025, in every negative quarter-hour, so such producers bid at zero or
stop. \textbf{14.} It buys (charges) in the cheapest hours, adding demand there, and sells
in the evening. \textbf{15.} Flexible consumers, storage, and producers that can stop and buy
back their forward sales cheaply. \textbf{16.} More export would have raised the German price
and lowered the neighbours', until the line filled or the prices met. \textbf{17.} Strong sun,
mild temperatures and holidays give high solar output against low demand. \textbf{18.} It is
paid to consume: its value is the negative price times the consumption it can shift.
\textbf{19.} \emph{Named result:} the unit should switch off if restarting costs less than
EUR~651\,308. \textbf{20.} That the system has more inflexible supply than demand at that
hour, and pays anyone who can absorb it or stop.
\end{solution}

\begin{solution}{iq:m3:power-markets-design:1}
By a daily sealed-bid auction at noon for each MTU of the next day: bids and offers form
stepwise curves, one algorithm clears all coupled zones together, maximising welfare subject
to cross-border capacities, and each zone gets a uniform price where its curves cross.
\iqlookfor{uniform price, welfare maximisation, \omterm{def:m3:power-markets-design:coupling}{market coupling}.}
\end{solution}

\begin{solution}{iq:m3:power-markets-design:2}
The supply stack sorted by marginal cost. Wind has near-zero marginal cost, so it shifts the
stack right and the demand line meets it lower: prices fall in windy hours, and more so
where the stack is steep.
\iqlookfor{shift of the stack, cannibalisation of windy hours.}
\end{solution}

\begin{solution}{iq:m3:power-markets-design:3}
Load, wind and solar forecasts (residual load), fuel and carbon prices, plant availability,
interconnector capacities and neighbours' residual load, calendar effects. In spring, solar
and wind forecasts dominate midday prices.
\iqlookfor{residual load as the key feature; weather.}
\end{solution}

\begin{solution}{iq:m3:power-markets-design:4}
When inflexible or subsidised supply exceeds demand plus exports: producers pay to avoid
restart costs or to keep subsidies. The auction's harmonised minimum bounds it (in Europe
$-\qty{500}{}$, then $-\qty{600}{EUR/MWh}$ from May 2026).
\iqlookfor{mechanisms, and knowledge of the floor.}
\end{solution}

\begin{solution}{iq:m3:power-markets-design:5}
Under zonal pricing it receives the zone's price whatever local congestion does, but may be
curtailed by redispatch; under nodal pricing it receives its node's price, which falls when
local lines are congested, so it bears the congestion risk and needs transmission rights to
hedge.
\iqlookfor{who bears congestion risk, and the hedge.}
\end{solution}

\begin{solution}{iq:m3:power-markets-design:6}
A welfare-maximisation problem with network constraints (a linear or quadratic programme
per scenario) plus a combinatorial layer for blocks (branch and bound, or a heuristic with
paradox checks). Speed: warm-start from the previous scenario, decompose by day, vectorise
the per-MTU curves, and cache block decisions that do not change.
\iqlookfor{LP plus combinatorics, and practical acceleration.}
\end{solution}
